\section*{Ethical Considerations}
\label{sec:ethics}

\textbf{Decision-support boundary.}
VIGIL emits risk intelligence, evidence summaries, and mitigation candidates with provenance and residual uncertainty for qualified human reviewers. It is not legal advice: it does not replace counsel, predict litigation outcomes, or recommend business decisions, and severity labels are triage proxies that are not interpreted as liability. Every High- or Medium-risk item is presented for review by a qualified person, and unresolved legal-status atoms are surfaced as such rather than resolved automatically.

\textbf{Human oversight of failure modes.}
Low-confidence negative-treatment edges enter a human review queue, since missed negative-treatment actions are the failure mode most likely to admit stale authority into a report. Reviewers can distinguish a genuine product--theory match from a retrieval gap or an unresolved legal-status ambiguity.

\textbf{Data governance.}
Released data are public docket metadata and machine-generated summaries; we do not redistribute attorney-authored complaint text beyond what is already public, and the P2 citator study releases only the crosswalk, sampling seeds, and agreement statistics, with no redistributed proprietary citator content. Source records are public U.S.\ court filings obtained through CourtListener/RECAP and are used for research consistent with that service's terms; news items are represented only by links and derived metadata. Released records contain the names of parties and counsel as they appear in these public filings; we collect no other personal information. VIGIL-Dataset is released under CC~BY~4.0 and the code under the MIT license, for research on legal-risk analysis; neither is intended for automated legal decisions. The Protocol~P2 comparison used a commercial citator under the authors' institutional research subscription; only agreement statistics and sampling seeds are released, with no redistributed citator content.

\textbf{Annotators and raters.}
The three legal research specialists who adjudicated labels and the two specialists who rated reports (\S\ref{sec:dataset}, \S\ref{sec:evaluation}) are staff of the authors' institution and performed this work within their regular paid duties, without additional compensation beyond ordinary salary. Before annotation began, participants received a written description of the task, the intended research use of their labels, data-handling practices, and the voluntary nature of participation; verbal consent was obtained and recorded in the project log. The work consisted solely of labeling public court documents and collected no personal information about the participants beyond institutional employment status. Under the authors' institutional human-subjects policy for secondary analysis of public records with no identifiable private information about research subjects, the authors determined that the activity does not constitute human-subjects research requiring IRB review (no separate IRB-office determination was sought); this determination is recorded in the released annotation guidelines. Institutional affiliation of annotators is a potential source of bias (Limitations). Annotation guidelines are released with the dataset so that external replications can use independent raters.

\textbf{Use of AI assistants.}
AI writing assistants were used to edit and polish the text of this paper. An AI coding assistant also wrote the code for the P1 pilot (Appendix~\ref{app:p1}), following a design fixed by the authors; the pilot design was frozen before any run, and its results were reported as produced. Research ideas, system design, and the remaining experiments and analyses are the authors' own.
